% Source: Evans, "Partial Differential Equations" (2nd ed., AMS Graduate Studies in Mathematics vol. 19, 2010),
% Chapter B "Inequalities", PDF pages 623-627 of MathBooks/Evans Partial Differential Equations(1).pdf.
% Transcribed page-by-page by vision subagents, then merged and restructured into
% leanblueprint conventions (semantic \label, bracketed titles, \uses dependency edges) below.
%
% NOTE ON GAPS: PDF page 626 was returned by the vision transcription subagent as a refusal
% placeholder rather than actual page content. That gap was closed by directly re-reading the
% corresponding page of the source PDF (book's own printed page 624) with a vision-capable Read
% call during this restructuring pass, so no content below is invented or reconstructed purely
% from outside knowledge of the book -- it is a faithful transcription of the actual page image.
%
% NOTE ON NOTATION: the raw OCR rendered the averaged-integral symbol as "\bar{\int}"; this has
% been normalized below to the project's established "\fint" notation (already used without
% definition in appendixA.tex and chapter5.tex, and supported natively by the KaTeX build).


\section{Convex Functions}
\label{sec:convex-functions}

\begin{definition}[Convex function]
\label{def:convex-function-rn}
\dcref{appB:B.1-def-1}
\group{Convex Analysis}
A function $f: \R^n \to \R$ is called \textit{convex} provided
$$(1) \qquad f(\tau x + (1 - \tau)y) \leq \tau f(x) + (1 - \tau)f(y)$$
for all $x, y \in \R^n$ and each $0 \leq \tau \leq 1$.
\end{definition}

\begin{theorem}[Supporting hyperplanes]
\label{thm:supporting-hyperplanes}
\dcref{appB:B.1-thm-1}
\group{Convex Analysis}
\uses{def:convex-function-rn}
Suppose $f: \R^n \to \R$ is \textit{convex}. Then for each $x \in \R^n$ there exists $r \in \R^n$ such that the inequality
$$(2) \qquad f(y) \geq f(x) + r \cdot (y - x)$$
holds for all $y \in \R^n$.
\end{theorem}

\begin{remark}[Supporting hyperplanes and uniform convexity]
\label{rem:supporting-hyperplane-remarks}
\dcref{appB:B.1-rem-1}
\group{Convex Analysis}
\uses{thm:supporting-hyperplanes}
(i) The mapping $y \mapsto f(x) + r \cdot (y - x)$ determines the \textit{supporting hyperplane} to $f$ at $x$. Inequality (2) says the graph of $f$ lies above each supporting hyperplane. If $f$ is differentiable at $x$, $r = Df(x)$.

(ii) If $f$ is $C^2$, then $f$ is \textit{convex} if and only if $D^2f \geq 0$. The $C^2$ function $f$ is \textit{uniformly convex} if $D^2f \geq \theta I$ for some constant $\theta > 0$: this means
$$\sum_{i,j=1}^{n} f_{x_i x_j}(x) \xi_i \xi_j \geq \theta |\xi|^2 \quad (x, \xi \in \R^n).$$
\end{remark}

\begin{theorem}[Jensen's inequality]
\label{thm:jensens-inequality}
\dcref{appB:B.1-thm-2}
\group{Convex Analysis}
\uses{def:convex-function-rn,thm:supporting-hyperplanes,def:function-notation-basic}
Assume $f: \R \to \R$ is \textit{convex}, and $U \subset \R^n$ is open, bounded. Let $u: U \to \R$ be summable. Then
$$(3) \qquad f\left(\fint_U u\,dx\right) \leq \fint_U f(u)\,dx.$$

Remember from \S A.3 the notation $\fint_U u\,dx = \frac{1}{|U|} \int_U u\,dx = $ average of $u$ over $U$.
\end{theorem}

\begin{proof}
Since $f$ is \textit{convex}, for each $p \in \R$ there exists $r \in \R$ such that
$$f(q) \geq f(p) + r(q - p) \quad \text{for all } q \in \R.$$
Let $p = \fint_U u\,dx$, $q = u(x)$:
$$f(u(x)) \geq f\left(\fint_U u\,dx\right) + r\left(u(x) - \fint_U u\,dx\right).$$
Integrate with respect to $x$ over $U$.
\end{proof}

Convex functions are discussed more fully in \S 3.3.2 and \S 9.6.1.

\section{Elementary Inequalities}
\label{sec:elementary-inequalities}

Following is a collection of elementary, but fundamental, inequalities. These estimates are continually employed throughout the text and should be memorized.

\begin{lemma}[Cauchy's inequality]
\label{lem:cauchys-inequality}
\dcref{appB:B.2-a}
\group{Elementary Inequalities}
$$(4) \qquad ab \leq \frac{a^2}{2} + \frac{b^2}{2} \quad (a,b \in \R).$$
\end{lemma}

\begin{proof}
$0 \leq (a-b)^2 = a^2 - 2ab + b^2$.
\end{proof}

\begin{lemma}[Cauchy's inequality with $\epsilon$]
\label{lem:cauchys-inequality-epsilon}
\dcref{appB:B.2-b}
\group{Elementary Inequalities}
\uses{lem:cauchys-inequality}
$$(5) \qquad ab \leq \epsilon a^2 + \frac{b^2}{4\epsilon} \quad (a,b > 0, \epsilon > 0).$$
\end{lemma}

\begin{proof}
Write
$$ab = ((2\epsilon)^{1/2}a)\left(\frac{b}{(2\epsilon)^{1/2}}\right)$$
and apply Cauchy's inequality.
\end{proof}

\begin{lemma}[Young's inequality]
\label{lem:youngs-inequality}
\dcref{appB:B.2-c}
\group{Elementary Inequalities}
\uses{def:convex-function-rn}
Let $1 < p, q < \infty$, $\frac{1}{p} + \frac{1}{q} = 1$. Then
$$(6) \qquad ab \leq \frac{a^p}{p} + \frac{b^q}{q} \quad (a,b > 0).$$
\end{lemma}

\begin{proof}
The mapping $x \mapsto e^x$ is convex, and consequently
$$ab = e^{\log a + \log b} = e^{\frac{1}{p}\log a^p + \frac{1}{q}\log b^q} \leq \frac{1}{p}e^{\log a^p} + \frac{1}{q}e^{\log b^q} = \frac{a^p}{p} + \frac{b^q}{q}.$$
\end{proof}

\begin{lemma}[Young's inequality with $\epsilon$]
\label{lem:youngs-inequality-epsilon}
\dcref{appB:B.2-d}
\group{Elementary Inequalities}
\uses{lem:youngs-inequality}
$$(7) \qquad ab \leq \epsilon a^p + C(\epsilon)b^q \quad (a,b > 0, \epsilon > 0)$$
for $C(\epsilon) = (\epsilon p)^{-q/p}q^{-1}$.
\end{lemma}

\begin{proof}
Write $ab = ((\epsilon p)^{1/p}a)\left(\frac{b}{(\epsilon p)^{1/p}}\right)$ and apply Young's inequality.
\end{proof}

\begin{lemma}[Hölder's inequality]
\label{lem:holders-inequality}
\dcref{appB:B.2-e}
\group{Elementary Inequalities}
\uses{lem:youngs-inequality}
Assume $1 \leq p, q \leq \infty$, $\frac{1}{p} + \frac{1}{q} = 1$. Then if $u \in L^p(U)$, $v \in L^q(U)$, we have
$$(8) \qquad \int_U |uv|\,dx \leq \|u\|_{L^p(U)}\|v\|_{L^q(U)}.$$
\end{lemma}

\begin{proof}
By homogeneity, we may assume $\|u\|_{L^p} = \|v\|_{L^q} = 1$. Then Young's inequality implies for $1 < p, q < \infty$ that
$$\int_U |uv|\,dx \leq \frac{1}{q}\int_U |u|^p\,dx + \frac{1}{q}\int_U |v|^q\,dx = 1 = \|u\|_{L^p}\|v\|_{L^q}.$$
\end{proof}

\begin{lemma}[Minkowski's inequality]
\label{lem:minkowskis-inequality}
\dcref{appB:B.2-f}
\group{Elementary Inequalities}
\uses{lem:holders-inequality}
Assume $1 \leq p \leq \infty$ and $u,v \in L^p(U)$. Then
$$(9) \qquad \|u + v\|_{L^p(U)} \leq \|u\|_{L^p(U)} + \|v\|_{L^p(U)}.$$
\end{lemma}

\begin{proof}
$$\|u + v\|_{L^p(U)}^p = \int_U |u + v|^p\,dx \leq \int_U |u + v|^{p-1}(|u| + |v|) \, dx$$
$$\leq \left(\int_U |u + v|^p\,dx\right)^{\frac{p-1}{p}} \left(\left(\int_U |u|^p\,dx\right)^{1/p} + \left(\int_U |v|^p\,dx\right)^{1/p}\right)$$
$$= \|u + v\|_{L^p(U)}^{p-1}(\|u\|_{L^p(U)} + \|v\|_{L^p(U)}).$$
\end{proof}

\begin{remark}[Discrete Hölder and Minkowski inequalities]
\label{rem:discrete-holder-minkowski-inequalities}
\dcref{appB:B.2-f-rem}
\group{Elementary Inequalities}
\uses{lem:holders-inequality,lem:minkowskis-inequality}
Similar proofs establish the discrete versions of Hölder's and Minkowski's inequalities:
$$(10) \qquad \begin{cases}
\left|\sum_{k=1}^n a_k b_k\right| \leq \left(\sum_{k=1}^n |a_k|^p\right)^{\frac{1}{p}} \left(\sum_{k=1}^n |b_k|^q\right)^{\frac{1}{q}}, \\
\left(\sum_{k=1}^n |a_k + b_k|^p\right)^{\frac{1}{p}} \leq \left(\sum_{k=1}^n |a_k|^p\right)^{\frac{1}{p}} + \left(\sum_{k=1}^n |b_k|^p\right)^{\frac{1}{p}},
\end{cases}$$
for $a = (a_1,\ldots,a_n), b = (b_1,\ldots,b_n) \in \R^n$ and $1 \leq p < \infty$, $\frac{1}{p} + \frac{1}{q} = 1$. $\square$
\end{remark}

\begin{lemma}[General Hölder inequality]
\label{lem:general-holder-inequality}
\dcref{appB:B.2-g}
\group{Elementary Inequalities}
\uses{lem:holders-inequality}
Let $1 \leq p_1,\ldots,p_m \leq \infty$, with $\frac{1}{p_1} + \frac{1}{p_2} + \cdots + \frac{1}{p_m} = 1$, and assume $u_k \in L^{p_k}(U)$ for $k = 1,\ldots,m$. Then
$$(11) \qquad \int_U |u_1 \cdots u_m| \, dx \leq \prod_{k=1}^m \|u_k\|_{L^{p_k}(U)}.$$
\end{lemma}

\begin{proof}
Induction, using Hölder's inequality.
\end{proof}

\begin{lemma}[Interpolation inequality for $L^p$-norms]
\label{lem:interpolation-inequality-lp-norms}
\dcref{appB:B.2-h}
\group{Elementary Inequalities}
\uses{lem:holders-inequality}
Assume $1 \leq s \leq r \leq t \leq \infty$ and
$$\frac{1}{r} = \frac{\theta}{s} + \frac{(1-\theta)}{t}.$$
Suppose also $u \in L^s(U) \cap L^t(U)$. Then $u \in L^r(U)$, and
$$(12) \qquad \|u\|_{L^r(U)} \leq \|u\|_{L^s(U)}^\theta \|u\|_{L^t(U)}^{1-\theta}.$$
\end{lemma}

\begin{proof}
We compute
$$\int_U |u|^r\,dx = \int_U |u|^{\theta r}|u|^{(1-\theta)r}\,dx$$
$$\leq \left(\int_U |u|^{\theta r \cdot \frac{s}{\theta r}}\,dx\right)^{\frac{\theta r}{s}} \left(\int_U |u|^{(1-\theta)r \cdot \frac{t}{(1-\theta) r}}\,dx\right)^{\frac{(1-\theta) r}{t}}.$$
We have invoked Hölder's inequality, which applies since $\frac{\theta r}{s} + \frac{(1-\theta)r}{t} = 1$.
\end{proof}

\begin{lemma}[Cauchy--Schwarz inequality]
\label{lem:cauchy-schwarz-inequality}
\dcref{appB:B.2-i}
\group{Elementary Inequalities}
$$(13) \qquad |x \cdot y| \leq |x||y| \quad (x,y \in \R^n).$$
\end{lemma}

\begin{proof}
Let $\epsilon > 0$ and note
$$0 \leq |x \pm \epsilon y|^2 = |x|^2 \pm 2\epsilon x \cdot y + \epsilon^2 |y|^2.$$
Consequently
$$\pm x \cdot y \leq \frac{1}{2\epsilon}|x|^2 + \frac{\epsilon}{2}|y|^2.$$
Minimize the right hand side by setting $\epsilon = \frac{|x|}{|y|}$, provided $y \neq 0$.
\end{proof}

\begin{remark}[Cauchy--Schwarz inequality for a symmetric nonnegative matrix]
\label{rem:cauchy-schwarz-matrix-form}
\dcref{appB:B.2-i-rem}
\group{Elementary Inequalities}
\uses{lem:cauchy-schwarz-inequality}
Likewise, if $A$ is a symmetric, nonnegative $n \times n$ matrix,
$$(14) \qquad \left|\sum_{i,j=1}^{n} a_{ij}x_iy_j\right| \leq \left(\sum_{i,j=1}^{n} a_{ij}x_ix_j\right)^{1/2} \left(\sum_{i,j=1}^{n} a_{ij}y_iy_j\right)^{1/2} \quad (x,y \in \R^n).$$
$\square$
\end{remark}

\begin{lemma}[Gronwall's inequality (differential form)]
\label{lem:gronwalls-inequality-differential-form}
\dcref{appB:B.2-j}
\group{Gronwall's Inequality}
(i) Let $\eta(\cdot)$ be a nonnegative, absolutely continuous function on $[0,T]$, which satisfies for a.e. $t$ the differential inequality
$$(15) \qquad \eta'(t) \leq \phi(t)\eta(t) + \psi(t),$$
where $\phi(t)$ and $\psi(t)$ are nonnegative, summable functions on $[0,T]$. Then
$$(16) \qquad \eta(t) \leq e^{\int_0^t \phi(s)\,ds}\left[\eta(0) + \int_0^t \psi(s)\,ds\right]$$
for all $0 \leq t \leq T$.

(ii) In particular, if
$$\eta' \leq \phi\eta \quad \text{on } [0,T] \quad \text{and} \quad \eta(0) = 0,$$
then
$$\eta \equiv 0 \quad \text{on } [0,T].$$
\end{lemma}

\begin{proof}
From (15) we see
$$\frac{d}{ds}\left(\eta(s)e^{-\int_0^s \phi(r)\,dr}\right) = e^{-\int_0^s \phi(r)\,dr}(\eta'(s) - \phi(s)\eta(s)) \le e^{-\int_0^s \phi(r)\,dr}\psi(s)$$
for a.e. $0 \le s \le T$. Consequently for each $0 \le t \le T$, we have
$$\eta(t)e^{-\int_0^t \phi(r)\,dr} \le \eta(0) + \int_0^t e^{-\int_0^s \phi(r)\,dr}\psi(s)\,ds \le \eta(0) + \int_0^t \psi(s)\,ds.$$
This implies inequality (16).
\end{proof}

\begin{lemma}[Gronwall's inequality (integral form)]
\label{lem:gronwalls-inequality-integral-form}
\dcref{appB:B.2-k}
\group{Gronwall's Inequality}
\uses{lem:gronwalls-inequality-differential-form}
(i) Let $\xi(t)$ be a nonnegative, summable function on $[0,T]$ which satisfies for a.e. $t$ the integral inequality
$$(17) \qquad \xi(t) \le C_1 \int_0^t \xi(s)\,ds + C_2$$
for constants $C_1, C_2 \ge 0$. Then
$$(18) \qquad \xi(t) \le C_2(1 + C_1 te^{C_1 t})$$
for a.e. $0 \le t \le T$.

(ii) In particular, if
$$\xi(t) \le C_1 \int_0^t \xi(s)\,ds$$
for a.e. $0 \le t \le T$, then
$$\xi(t) = 0 \quad \text{a.e.}$$
\end{lemma}

\begin{proof}
Let $\eta(t) := \int_0^t \xi(s)\,ds$; then $\eta \le C_1 \eta + C_2$ a.e. in $[0,T]$. According to the differential form of Gronwall's inequality above:
$$\eta(t) \le e^{C_1 t}(\eta(0) + C_2 t) = C_2 te^{C_1 t}.$$
Then (17) implies
$$\xi(t) \le C_1 \eta(t) + C_2 \le C_2(1 + C_1 te^{C_1 t}).$$
\end{proof}
